% =========================================================
% Chapter: Ideal Gases, Kinetic Theory, and the First Law of Thermodynamics
% POSN 1 Physics Preparation Notes
% Language: Thai
% Compiler: XeLaTeX
% =========================================================

\chapter{แก๊สอุดมคติ ทฤษฎีจลน์ และกฎข้อที่หนึ่งของอุณหพลศาสตร์}

\begin{center}
    {\Large \textbf{Ideal Gases, Kinetic Theory, and the First Law of Thermodynamics}}\\[4pt]
    {\normalsize เอกสารเตรียมสอบคัดเลือก สอวน. ฟิสิกส์ ค่าย 1}\\[6pt]
    {\normalsize จัดทำโดย \textbf{Tames Thanakrit Damduan}}\\[2pt]
    {\footnotesize \textcopyright\ 2026 Tames Thanakrit Damduan. All rights reserved.}\\[8pt]
\end{center}

\section*{เป้าหมายของบทเรียน}

หลังจากเรียนบทนี้ นักเรียนควรทำสิ่งต่อไปนี้ได้

\begin{enumerate}
    \item ใช้กฎของแก๊สอุดมคติในรูป $PV=nRT$ และ $PV=Nk_BT$ ได้
    \item แปลงอุณหภูมิเป็นเคลวินก่อนใช้สมการของแก๊สได้อย่างถูกต้อง
    \item แยกความดันสัมบูรณ์ออกจากความดันเกจ และใช้ความดันที่ถูกต้องในกฎของแก๊สได้
    \item ใช้ความสัมพันธ์ระหว่างสถานะเริ่มต้นและสถานะสุดท้ายของแก๊สปริมาณคงตัวได้
    \item เชื่อมโยงความดันของแก๊สกับการชนของโมเลกุลตามแบบจำลองทฤษฎีจลน์ได้
    \item ใช้ความสัมพันธ์ระหว่างอุณหภูมิกับพลังงานจลน์เฉลี่ยและอัตราเร็วรากกำลังสองเฉลี่ยได้
    \item วิเคราะห์ลักษณะของกราฟการกระจายอัตราเร็วของโมเลกุลได้
    \item ใช้หลักการแบ่งเท่าเทียมของพลังงานและหาพลังงานภายในของแก๊สอุดมคติอะตอมเดี่ยวได้
    \item คำนวณงานที่แก๊สทำจากกราฟ $P$ กับ $V$ และใช้กฎข้อที่หนึ่งของอุณหพลศาสตร์ได้
    \item แก้โจทย์ประยุกต์ระดับ สอวน. ที่เชื่อมกฎของแก๊สกับของไหล แรงพยุง และการประมาณค่าได้
\end{enumerate}

\begin{tcolorbox}[title=ขอบเขตของบทนี้]
บทก่อนหน้าได้กล่าวถึงอุณหภูมิ สมดุลความร้อน ความร้อนแฝง การขยายตัวเชิงความร้อน และการนำความร้อนแล้ว

บทนี้ศึกษาพฤติกรรมของแก๊สตั้งแต่ระดับมหภาคไปสู่ระดับจุลภาค และจบด้วยกฎข้อที่หนึ่งของอุณหพลศาสตร์ เนื้อหาที่เน้นคือหัวข้อที่สามารถใช้ในข้อสอบคัดเลือกค่าย 1 ได้แก่

\[
PV=nRT
\]

\[
\frac{1}{2}k_BT \text{ ต่อหนึ่ง degree of freedom}
\]

\[
\Delta U=Q-W
\]

โดย $W$ หมายถึงงานที่แก๊สทำต่อสิ่งแวดล้อม
\end{tcolorbox}

\newpage{}

% =========================================================
\section{ตัวแปรสถานะของแก๊ส}

\subsection{ความดัน ปริมาตร และอุณหภูมิ}

สถานะของแก๊สในภาวะสมดุลสามารถบรรยายด้วยปริมาณสำคัญ ได้แก่

\[
P=\text{ความดัน}
\]

\[
V=\text{ปริมาตร}
\]

\[
T=\text{อุณหภูมิสัมบูรณ์}
\]

\[
n=\text{จำนวนโมล}
\]

ในการใช้สมการของแก๊ส ต้องใช้อุณหภูมิหน่วยเคลวินเสมอ

\[
\boxed{T_{\si{K}}=T_{\si{\degreeCelsius}}+273.15}
\]

สำหรับข้อสอบที่ไม่ต้องการความละเอียดสูง มักใช้

\[
T_{\si{K}}\approx T_{\si{\degreeCelsius}}+273
\]

\begin{tcolorbox}[title=ข้อควรระวัง]
ในสมการแคลอริเมตรี $Q=mc\Delta T$ สามารถใช้ผลต่างอุณหภูมิในหน่วยองศาเซลเซียสได้

แต่ในสมการของแก๊ส

\[
PV=nRT
\]

ต้องใช้ค่าอุณหภูมิสัมบูรณ์ในหน่วยเคลวินเท่านั้น
\end{tcolorbox}

\subsection{จำนวนโมลและจำนวนโมเลกุล}

ถ้าแก๊สมีจำนวนโมเลกุลทั้งหมด $N$ และมีจำนวนโมล $n$

\[
N=nN_A
\]

โดย

\[
N_A=6.02\times10^{23}\,\si{mol^{-1}}
\]

คือเลขอาโวกาโดร

ถ้าโมเลกุลหนึ่งมีมวล $m_0$ และแก๊สมีมวลรวม $m$

\[
m=Nm_0
\]

ถ้าใช้มวลโมลาร์ $M_{\text{mol}}$

\[
m=nM_{\text{mol}}
\]

% =========================================================
\section{กฎของแก๊สอุดมคติ}

\subsection{สมการระดับมหภาค}

แก๊สอุดมคติเป็นแบบจำลองที่ใช้อธิบายแก๊สจริงได้ดีเมื่อความหนาแน่นไม่สูงมาก และแรงยึดเหนี่ยวระหว่างโมเลกุลไม่มีบทบาทเด่น

สมการของแก๊สอุดมคติคือ

\[
\boxed{PV=nRT}
\]

โดย

\[
R=8.31\,\si{J/(mol.K)}
\]

ถ้าใช้จำนวนโมเลกุลแทนจำนวนโมล จะเขียนได้เป็น

\[
\boxed{PV=Nk_BT}
\]

โดย

\[
k_B=1.38\times10^{-23}\,\si{J/K}
\]

คือค่าคงที่โบลทซ์มันน์

เพราะ

\[
N=nN_A
\]

จึงได้ความสัมพันธ์

\[
\boxed{R=N_Ak_B}
\]

\begin{examplebox}{ตัวอย่าง: หาจำนวนโมเลกุลของแก๊ส}
แก๊สอุดมคติอยู่ในภาชนะปริมาตร $0.050\,\si{m^3}$ ที่ความดัน $1.0\times10^5\,\si{Pa}$ และอุณหภูมิ $300\,\si{K}$ จงหาจำนวนโมเลกุลโดยประมาณ

ใช้

\[
PV=Nk_BT
\]

ดังนั้น

\[
N=\frac{PV}{k_BT}
\]

แทนค่า

\[
N=\frac{(1.0\times10^5)(0.050)}{(1.38\times10^{-23})(300)}
\]

\[
\boxed{N\approx1.2\times10^{24}\text{ โมเลกุล}}
\]
\end{examplebox}

\subsection{แก๊สปริมาณคงตัวเปลี่ยนสถานะ}

ถ้าจำนวนโมลของแก๊สไม่เปลี่ยน

\[
PV=nRT
\]

ทำให้

\[
\frac{PV}{T}=nR=\text{ค่าคงตัว}
\]

ดังนั้นระหว่างสถานะเริ่มต้นและสถานะสุดท้าย

\[
\boxed{
\frac{P_1V_1}{T_1}
=
\frac{P_2V_2}{T_2}
}
\]

\subsection{กรณีพิเศษที่ควรรู้}

\begin{enumerate}
    \item อุณหภูมิคงตัว
    \[
    T=\text{ค่าคงตัว}
    \]
    จึงได้
    \[
    \boxed{P_1V_1=P_2V_2}
    \]

    \item ความดันคงตัว
    \[
    P=\text{ค่าคงตัว}
    \]
    จึงได้
    \[
    \boxed{\frac{V_1}{T_1}=\frac{V_2}{T_2}}
    \]

    \item ปริมาตรคงตัว
    \[
    V=\text{ค่าคงตัว}
    \]
    จึงได้
    \[
    \boxed{\frac{P_1}{T_1}=\frac{P_2}{T_2}}
    \]
\end{enumerate}

\subsection{ความดันสัมบูรณ์และความดันเกจ}

กฎของแก๊สอุดมคติต้องใช้ความดันสัมบูรณ์ ซึ่งวัดเทียบกับสุญญากาศ

ถ้าโจทย์ให้ความดันเกจ

\[
P_{\text{gauge}}
\]

ต้องแปลงเป็นความดันสัมบูรณ์ก่อน

\[
\boxed{
P_{\text{abs}}
=
P_{\text{atm}}
+
P_{\text{gauge}}
}
\]

ถ้าความดันภายในต่ำกว่าความดันบรรยากาศ สามารถเขียนเป็น

\[
P_{\text{abs}}
=
P_{\text{atm}}-\Delta P
\]

เมื่อ $\Delta P$ คือขนาดผลต่างความดัน

\begin{examplebox}{ตัวอย่าง: บอลลูนปริมาตรคงตัว}
แก๊สในบอลลูนปริมาตรคงตัวมีความดันเกจ $200\,\si{kPa}$ ที่ระดับพื้นดิน ซึ่งมีความดันบรรยากาศ $100\,\si{kPa}$ และอุณหภูมิ $300\,\si{K}$ เมื่อบอลลูนลอยขึ้นไปยังบริเวณที่อุณหภูมิ $240\,\si{K}$ และความดันบรรยากาศ $60\,\si{kPa}$ จงหาความดันเกจใหม่

ความดันสัมบูรณ์เริ่มต้นคือ

\[
P_1=100+200=300\,\si{kPa}
\]

เพราะปริมาตรคงตัว

\[
\frac{P_1}{T_1}=\frac{P_2}{T_2}
\]

\[
P_2=300\left(\frac{240}{300}\right)
\]

\[
P_2=240\,\si{kPa}
\]

ดังนั้นความดันเกจใหม่คือ

\[
P_{\text{gauge},2}=240-60
\]

\[
\boxed{P_{\text{gauge},2}=180\,\si{kPa}}
\]
\end{examplebox}

\subsection{ความหนาแน่นของแก๊สอุดมคติ}

ถ้าแก๊สมีมวลโมลาร์ $M_{\text{mol}}$

\[
m=nM_{\text{mol}}
\]

ความหนาแน่นคือ

\[
\rho=\frac{m}{V}
\]

จาก

\[
PV=nRT
\]

จะได้

\[
\frac{n}{V}=\frac{P}{RT}
\]

ดังนั้น

\[
\boxed{
\rho
=
\frac{PM_{\text{mol}}}{RT}
}
\]

ที่ความดันและอุณหภูมิเดียวกัน ความหนาแน่นจึงแปรผันตรงกับมวลโมลาร์

\[
\boxed{
ho\propto M_{\text{mol}}}
\]

\subsection{การเปลี่ยนแปลงเล็กน้อยของแก๊ส}

ถ้าแก๊สมีอุณหภูมิคงตัว

\[
PV=\text{ค่าคงตัว}
\]

เมื่อ $P$ และ $V$ เปลี่ยนไปเล็กน้อย

\[
(P+\Delta P)(V+\Delta V)=PV
\]

ละพจน์อันดับสอง $\Delta P\Delta V$ จะได้

\[
P\Delta V+V\Delta P\approx0
\]

ดังนั้น

\[
\boxed{
\frac{\Delta P}{\Delta V}
\approx
-\frac{P}{V}
}
\]

ในภาษาของแคลคูลัส

\[
\boxed{
\left(\frac{dP}{dV}\right)_T
=
-\frac{P}{V}
}
\]

\begin{exercisebox}{แบบฝึกหัด 1: กฎของแก๊สอุดมคติ}
\begin{enumerate}
    \item แก๊สอุดมคติ $2.0\,\si{mol}$ อยู่ที่อุณหภูมิ $300\,\si{K}$ และความดัน $1.0\times10^5\,\si{Pa}$ จงหาปริมาตร
    \item แก๊สในถังแข็งมีอุณหภูมิเพิ่มจาก $300\,\si{K}$ เป็น $450\,\si{K}$ ถ้าความดันเริ่มต้นคือ $80\,\si{kPa}$ จงหาความดันสุดท้าย
    \item ที่ความดันและอุณหภูมิเดียวกัน แก๊สฮีเลียมและแก๊สออกซิเจนมีความหนาแน่นต่างกันกี่เท่า กำหนดให้มวลโมลาร์เป็น $4$ และ $32\,\si{g/mol}$ ตามลำดับ
    \item แก๊สปริมาณคงตัวอยู่ในสภาวะอุณหภูมิคงตัว ถ้าปริมาตรเพิ่มขึ้นเล็กน้อย $1\%$ ความดันจะเปลี่ยนไปประมาณกี่เปอร์เซ็นต์
\end{enumerate}
\end{exercisebox}

% =========================================================
\section{แบบจำลองจุลภาคของแก๊สอุดมคติ}

\subsection{สมมติฐานพื้นฐานของทฤษฎีจลน์}

ทฤษฎีจลน์อธิบายสมบัติระดับมหภาคของแก๊สจากการเคลื่อนที่ของโมเลกุลจำนวนมาก แบบจำลองพื้นฐานมีสมมติฐานดังนี้

\begin{enumerate}
    \item โมเลกุลเคลื่อนที่แบบสุ่มตลอดเวลา
    \item ขนาดของโมเลกุลเล็กมากเมื่อเทียบกับระยะห่างเฉลี่ยระหว่างโมเลกุล
    \item แรงระหว่างโมเลกุลละได้ ยกเว้นในช่วงเวลาสั้น ๆ ระหว่างการชน
    \item การชนระหว่างโมเลกุล และการชนกับผนัง เป็นการชนยืดหยุ่น
    \item กฎของนิวตันใช้อธิบายการเคลื่อนที่ของโมเลกุลแต่ละตัวได้
\end{enumerate}

\subsection{ความดันเกิดจากการชนกับผนัง}

ให้แก๊สมีโมเลกุลทั้งหมด $N$ แต่ละโมเลกุลมีมวล $m_0$ อยู่ในภาชนะปริมาตร $V$

เมื่อโมเลกุลชนผนัง โมเมนตัมตั้งฉากกับผนังเปลี่ยนไป แรงเฉลี่ยจากการชนของโมเลกุลจำนวนมากทำให้เกิดความดัน

ผลจากการวิเคราะห์ในสามมิติคือ

\[
\boxed{
P
=
\frac{1}{3}
\left(\frac{N}{V}\right)
 m_0\overline{v^2}
}
\]

หรือ

\[
\boxed{
PV
=
\frac{1}{3}Nm_0\overline{v^2}
}
\]

โดย

\[
\overline{v^2}
\]

คือค่าเฉลี่ยของกำลังสองของอัตราเร็ว ไม่ใช่กำลังสองของอัตราเร็วเฉลี่ย

\subsection{อัตราเร็วรากกำลังสองเฉลี่ย}

นิยาม

\[
\boxed{
v_{\text{rms}}
=
\sqrt{\overline{v^2}}
}
\]

จาก

\[
PV
=
\frac{1}{3}Nm_0\overline{v^2}
\]

และ

\[
PV=Nk_BT
\]

จะได้

\[
Nk_BT
=
\frac{1}{3}Nm_0\overline{v^2}
\]

ตัด $N$

\[
k_BT
=
\frac{1}{3}m_0\overline{v^2}
\]

ดังนั้น

\[
\boxed{
\frac{1}{2}m_0\overline{v^2}
=
\frac{3}{2}k_BT
}
\]

และ

\[
\boxed{
 v_{\text{rms}}
=
\sqrt{\frac{3k_BT}{m_0}}
}
\]

ถ้าใช้มวลโมลาร์ $M_{\text{mol}}$

\[
\boxed{
 v_{\text{rms}}
=
\sqrt{\frac{3RT}{M_{\text{mol}}}}
}
\]

โดยต้องใช้ $M_{\text{mol}}$ ในหน่วย $\si{kg/mol}$

\begin{tcolorbox}[title=ความหมายของอุณหภูมิในทฤษฎีจลน์]
สำหรับแก๊สอุดมคติ อุณหภูมิสัมบูรณ์เป็นตัววัดพลังงานจลน์เฉลี่ยของการเคลื่อนที่แบบเลื่อนที่ของโมเลกุล

\[
\left\langle K_{\text{trans}}\right\rangle
=
\frac{3}{2}k_BT
\]

ไม่ได้หมายความว่าโมเลกุลทุกตัวมีอัตราเร็วเท่ากัน
\end{tcolorbox}

\begin{examplebox}{ตัวอย่าง: เปรียบเทียบอัตราเร็วของโมเลกุล}
แก๊สสองชนิดมีอุณหภูมิเท่ากัน โมเลกุลของแก๊ส $A$ มีมวลเป็น $4$ เท่าของโมเลกุลแก๊ส $B$ จงหาอัตราส่วน $v_{\text{rms},A}/v_{\text{rms},B}$

เพราะ

\[
v_{\text{rms}}\propto\frac{1}{\sqrt{m_0}}
\]

จึงได้

\[
\frac{v_{\text{rms},A}}{v_{\text{rms},B}}
=
\sqrt{\frac{m_B}{m_A}}
\]

\[
\frac{v_{\text{rms},A}}{v_{\text{rms},B}}
=
\sqrt{\frac{1}{4}}
\]

\[
\boxed{
\frac{v_{\text{rms},A}}{v_{\text{rms},B}}
=
\frac{1}{2}
}
\]
\end{examplebox}

% =========================================================
\section{การกระจายอัตราเร็วของโมเลกุล}

\subsection{โมเลกุลไม่ได้มีอัตราเร็วเท่ากันทั้งหมด}

แม้แก๊สจะมีอุณหภูมิคงตัว โมเลกุลแต่ละตัวก็ยังชนกันตลอดเวลา ทำให้อัตราเร็วของแต่ละโมเลกุลเปลี่ยนไป

จำนวนโมเลกุลที่มีอัตราเร็วอยู่ใกล้ค่า $v$ สามารถแสดงด้วยกราฟการกระจายอัตราเร็ว

สำหรับโจทย์ระดับค่าย 1 ควรอ่านลักษณะของฟังก์ชันรูป

\[
N(v)
=
Av^2e^{-Bv^2/T}
\]

โดย $A$ และ $B$ เป็นค่าคงตัวบวก

\subsection{การอ่านรูปกราฟจากสมการ}

ที่

\[
v=0
\]

จะได้

\[
N(0)=0
\]

เมื่อ $v$ มีค่ามากมาก

\[
e^{-Bv^2/T}\longrightarrow0
\]

ดังนั้น

\[
N(v)\longrightarrow0
\]

ฟังก์ชันจึงเริ่มจากศูนย์ เพิ่มขึ้นจนถึงค่าสูงสุดหนึ่งครั้ง แล้วลดลงเข้าใกล้ศูนย์

ถ้าหาอนุพันธ์

\[
\frac{dN}{dv}
=
2Ave^{-Bv^2/T}
\left(1-\frac{Bv^2}{T}\right)
\]

ที่ยอดกราฟซึ่ง $v>0$

\[
1-\frac{Bv^2}{T}=0
\]

ดังนั้นอัตราเร็วที่ยอดกราฟคือ

\[
\boxed{
v_{\text{peak}}
=
\sqrt{\frac{T}{B}}
}
\]

เมื่ออุณหภูมิเพิ่มขึ้น ยอดกราฟจึงเลื่อนไปทางขวา

\begin{tcolorbox}[title=ลักษณะกราฟที่ต้องจำ]
กราฟการกระจายอัตราเร็วของโมเลกุลต้องมีลักษณะสำคัญดังนี้

\begin{enumerate}
    \item เริ่มจากศูนย์เมื่อ $v=0$
    \item เพิ่มขึ้นอย่างราบเรียบ
    \item มีค่าสูงสุดหนึ่งครั้ง
    \item ลดลงและเข้าใกล้ศูนย์เมื่อ $v$ มาก
    \item เมื่อ $T$ สูงขึ้น กราฟกว้างขึ้นและยอดเลื่อนไปทางขวา
\end{enumerate}
\end{tcolorbox}

% =========================================================
\section{หลักการแบ่งเท่าเทียมของพลังงาน}

\subsection{หนึ่ง degree of freedom ได้พลังงานเฉลี่ยเท่าไร}

หลักการแบ่งเท่าเทียมของพลังงานกล่าวว่า แต่ละพจน์พลังงานที่อยู่ในรูปกำลังสองมีพลังงานเฉลี่ย

\[
\boxed{
\frac{1}{2}k_BT
}
\]

ต่อหนึ่งโมเลกุล ต่อหนึ่ง degree of freedom

หรือ

\[
\boxed{
\frac{1}{2}RT
}
\]

ต่อหนึ่งโมล ต่อหนึ่ง degree of freedom

\subsection{แก๊สอุดมคติอะตอมเดี่ยว}

โมเลกุลของแก๊สอะตอมเดี่ยวเคลื่อนที่ได้สามทิศทาง

\[
x,
\qquad
y,
\qquad
z
\]

จึงมี degree of freedom สำหรับการเลื่อนที่ทั้งหมด

\[
f=3
\]

พลังงานจลน์เฉลี่ยต่อหนึ่งโมเลกุลคือ

\[
\left\langle K\right\rangle
=
3\left(\frac{1}{2}k_BT\right)
\]

\[
\boxed{
\left\langle K\right\rangle
=
\frac{3}{2}k_BT
}
\]

ถ้ามี $N$ โมเลกุล พลังงานภายในคือ

\[
U
=
N\left(\frac{3}{2}k_BT\right)
\]

เพราะ

\[
Nk_B=nR
\]

จึงได้

\[
\boxed{
U
=
\frac{3}{2}nRT
}
\]

และ

\[
\boxed{
\Delta U
=
\frac{3}{2}nR\Delta T
}
\]

\subsection{รูปทั่วไปเมื่อมี $f$ degrees of freedom}

ถ้าโจทย์กำหนดจำนวน degree of freedom เป็น $f$

\[
\boxed{
U
=
\frac{f}{2}nRT
}
\]

และ

\[
\boxed{
\Delta U
=
\frac{f}{2}nR\Delta T
}
\]

\begin{tcolorbox}[title=ขอบเขตสำหรับข้อสอบคัดเลือก]
ควรเข้าใจหลักทั่วไปว่าแต่ละ degree of freedom ให้พลังงานเฉลี่ย $\frac{1}{2}k_BT$

แต่โจทย์คำนวณที่พบบ่อยที่สุดคือแก๊สอุดมคติอะตอมเดี่ยว ซึ่งใช้

\[
U=\frac{3}{2}nRT
\]
\end{tcolorbox}

% =========================================================
\section{งานที่แก๊สทำ}

\subsection{ลูกสูบและงานจากการขยายตัว}

พิจารณาแก๊สในกระบอกสูบที่มีลูกสูบเคลื่อนที่ได้ ถ้าแก๊สดันลูกสูบออกไปเล็กน้อย ปริมาตรของแก๊สเพิ่มขึ้น $dV$

งานเล็กน้อยที่แก๊สทำคือ

\[
\boxed{dW=P\,dV}
\]

ดังนั้นงานรวมคือ

\[
\boxed{
W
=
\int_{V_i}^{V_f}P\,dV
}
\]

\subsection{เครื่องหมายของงาน}

\begin{enumerate}
    \item ถ้าแก๊สขยายตัว
    \[
    V_f>V_i
    \]
    จะได้
    \[
    W>0
    \]
    แก๊สทำงานต่อสิ่งแวดล้อม

    \item ถ้าแก๊สถูกอัด
    \[
    V_f<V_i
    \]
    จะได้
    \[
    W<0
    \]
    สิ่งแวดล้อมทำงานต่อแก๊ส
\end{enumerate}

\subsection{พื้นที่ใต้กราฟ $P$ กับ $V$}

บนกราฟที่แกนนอนเป็น $V$ และแกนตั้งเป็น $P$

\[
W
=
\int P\,dV
\]

จึงเท่ากับพื้นที่ใต้กราฟระหว่าง $V_i$ ถึง $V_f$

\begin{tcolorbox}[title=งานจากกราฟ $P$ กับ $V$]
\[
\boxed{
W=\text{พื้นที่ใต้กราฟ }P\text{ กับ }V
}
\]

ต้องใส่เครื่องหมายตามทิศทางของกระบวนการ
\end{tcolorbox}

\subsection{ความดันคงตัว}

ถ้า

\[
P=\text{ค่าคงตัว}
\]

จะได้

\[
\boxed{
W
=P(V_f-V_i)
=P\Delta V
}
\]

สำหรับแก๊สอุดมคติที่ความดันคงตัว

\[
P\Delta V=nR\Delta T
\]

ดังนั้น

\[
\boxed{W=nR\Delta T}
\]

\subsection{ปริมาตรคงตัว}

ถ้า

\[
V=\text{ค่าคงตัว}
\]

จะมี

\[
dV=0
\]

ดังนั้น

\[
\boxed{W=0}
\]

\subsection{อ่านเสริม: อุณหภูมิคงตัวของแก๊สอุดมคติ}

ถ้าอุณหภูมิคงตัว

\[
PV=nRT
\]

จึงมี

\[
P=\frac{nRT}{V}
\]

ดังนั้น

\[
W
=
\int_{V_i}^{V_f}\frac{nRT}{V}\,dV
\]

\[
\boxed{
W
=
nRT\ln\left(\frac{V_f}{V_i}\right)
}
\]

% =========================================================
\section{กฎข้อที่หนึ่งของอุณหพลศาสตร์}

\subsection{หลักอนุรักษ์พลังงานสำหรับระบบเชิงความร้อน}

กฎข้อที่หนึ่งของอุณหพลศาสตร์คือหลักอนุรักษ์พลังงานสำหรับระบบที่รับหรือคายพลังงานความร้อน และอาจทำงานต่อสิ่งแวดล้อม

ในบทนี้ใช้ข้อตกลงว่า

\[
Q>0
\]

เมื่อความร้อนไหลเข้าสู่ระบบ และ

\[
W>0
\]

เมื่อระบบหรือแก๊สทำงานต่อสิ่งแวดล้อม

ดังนั้น

\[
\boxed{
\Delta U
=
Q-W
}
\]

หรือเขียนเป็น

\[
\boxed{
Q
=
\Delta U+W
}
\]

\begin{tcolorbox}[title=ความหมายเชิงพลังงาน]
พลังงานความร้อนที่ใส่เข้าไปในระบบ แบ่งได้เป็นสองส่วน

\[
Q
=
\Delta U+W
\]

ส่วนหนึ่งเพิ่มพลังงานภายใน

\[
\Delta U
\]

และอีกส่วนถูกส่งออกไปเป็นงานที่แก๊สทำ

\[
W
\]
\end{tcolorbox}

\subsection{กระบวนการปริมาตรคงตัว}

ถ้าปริมาตรคงตัว

\[
W=0
\]

ดังนั้น

\[
\boxed{
\Delta U=Q
}
\]

สำหรับแก๊สอุดมคติอะตอมเดี่ยว

\[
\boxed{
Q
=
\frac{3}{2}nR\Delta T
}
\]

\subsection{กระบวนการความดันคงตัว}

สำหรับแก๊สอุดมคติอะตอมเดี่ยว

\[
\Delta U
=
\frac{3}{2}nR\Delta T
\]

และ

\[
W
=
P\Delta V
=
nR\Delta T
\]

ดังนั้น

\[
Q
=
\Delta U+W
\]

\[
Q
=
\frac{3}{2}nR\Delta T+nR\Delta T
\]

\[
\boxed{
Q
=
\frac{5}{2}nR\Delta T
}
\]

จึงได้

\[
\boxed{
\Delta T
=
\frac{2Q}{5nR}
}
\]

\subsection{กระบวนการอุณหภูมิคงตัว}

สำหรับแก๊สอุดมคติ พลังงานภายในขึ้นกับอุณหภูมิเท่านั้น

ถ้า

\[
\Delta T=0
\]

จะได้

\[
\Delta U=0
\]

จากกฎข้อที่หนึ่ง

\[
0=Q-W
\]

ดังนั้น

\[
\boxed{Q=W}
\]

\subsection{กระบวนการอะเดียแบติก}

ถ้าไม่มีความร้อนไหลเข้าหรือออกจากระบบ

\[
Q=0
\]

จาก

\[
\Delta U=Q-W
\]

จะได้

\[
\boxed{
\Delta U=-W
}
\]

ถ้าแก๊สขยายตัวและทำงาน

\[
W>0
\]

พลังงานภายในจะลดลง และอุณหภูมิมักลดลง

\subsection{กระบวนการเป็นวัฏจักร}

ถ้าระบบกลับมายังสถานะเดิม

\[
\Delta U=0
\]

ดังนั้น

\[
\boxed{Q=W}
\]

สำหรับหนึ่งรอบของวัฏจักร

\begin{examplebox}{ตัวอย่าง: ให้ความร้อนแก่แก๊สอะตอมเดี่ยวที่ความดันคงตัว}
ให้ความร้อน $Q$ แก่แก๊สอุดมคติอะตอมเดี่ยวจำนวน $n$ โมลอย่างช้า ๆ ในกระบอกสูบที่มีลูกสูบลื่น และความดันภายนอกคงตัว จงหาอุณหภูมิที่เพิ่มขึ้น

พลังงานภายในเพิ่มขึ้น

\[
\Delta U
=
\frac{3}{2}nR\Delta T
\]

งานที่แก๊สทำคือ

\[
W=P\Delta V
\]

เพราะความดันคงตัว

\[
P\Delta V=nR\Delta T
\]

ดังนั้น

\[
W=nR\Delta T
\]

ใช้

\[
Q=\Delta U+W
\]

\[
Q
=
\frac{3}{2}nR\Delta T+nR\Delta T
\]

\[
Q
=
\frac{5}{2}nR\Delta T
\]

ดังนั้น

\[
\boxed{
\Delta T
=
\frac{2Q}{5nR}
}
\]
\end{examplebox}

\begin{exercisebox}{แบบฝึกหัด 2: ทฤษฎีจลน์และกฎข้อที่หนึ่ง}
\begin{enumerate}
    \item แก๊สอุดมคติอะตอมเดี่ยว $2.0\,\si{mol}$ มีอุณหภูมิเพิ่มขึ้น $50\,\si{K}$ จงหาการเพิ่มของพลังงานภายใน
    \item แก๊สอุดมคติอะตอมเดี่ยวจำนวน $n$ โมลรับความร้อน $Q$ ที่ปริมาตรคงตัว จงหา $\Delta T$
    \item แก๊สอุดมคติอะตอมเดี่ยวจำนวน $n$ โมลรับความร้อน $Q$ ที่ความดันคงตัว จงหา $\Delta T$
    \item แก๊สอุดมคติขยายตัวแบบอุณหภูมิคงตัวจาก $V$ เป็น $3V$ ถ้าอุณหภูมิเป็น $T$ จงหางานที่แก๊สทำ
    \item ในกราฟ $P$ กับ $V$ แก๊สขยายตัวที่ความดันคงตัว $P_0$ จาก $V_0$ เป็น $4V_0$ แล้วลดความดันที่ปริมาตรคงตัวลงเหลือ $P_0/2$ จงหางานรวมที่แก๊สทำ
    \item โมเลกุลแก๊สชนิดหนึ่งมีมวลเป็น $9$ เท่าของอีกชนิดหนึ่ง ถ้าอุณหภูมิเท่ากัน จงหาอัตราส่วนของ $v_{\text{rms}}$
\end{enumerate}
\end{exercisebox}

% =========================================================
\section{วิธีเลือกหลักในการแก้โจทย์}

\begin{tcolorbox}[title=แผนการเลือกสมการ]
\begin{enumerate}
    \item ถ้าโจทย์ให้ $P$, $V$, $T$, จำนวนโมล หรือจำนวนโมเลกุล ใช้
    \[
    PV=nRT
    \]
    หรือ
    \[
    PV=Nk_BT
    \]

    \item ถ้าแก๊สปริมาณเดิมเปลี่ยนสถานะ ใช้
    \[
    \frac{P_1V_1}{T_1}
    =
    \frac{P_2V_2}{T_2}
    \]

    \item ถ้าโจทย์ให้ความดันเกจ ต้องเปลี่ยนเป็นความดันสัมบูรณ์ก่อน
    \[
    P_{\text{abs}}
    =
    P_{\text{atm}}+P_{\text{gauge}}
    \]

    \item ถ้าโจทย์ถามความหนาแน่นของแก๊ส ใช้
    \[
    \rho
    =
    \frac{PM_{\text{mol}}}{RT}
    \]

    \item ถ้าโจทย์ถามอัตราเร็วระดับโมเลกุล ใช้
    \[
    v_{\text{rms}}
    =
    \sqrt{\frac{3k_BT}{m_0}}
    \]

    \item ถ้าโจทย์ถามพลังงานภายในของแก๊สอะตอมเดี่ยว ใช้
    \[
    U
    =
    \frac{3}{2}nRT
    \]

    \item ถ้าโจทย์ถามงานจากกราฟ $P$ กับ $V$ ใช้พื้นที่ใต้กราฟ หรือ
    \[
    W=\int P\,dV
    \]

    \item ถ้าโจทย์มีความร้อน งาน และพลังงานภายใน ใช้
    \[
    \Delta U=Q-W
    \]
\end{enumerate}
\end{tcolorbox}

% =========================================================
\section{ชุดโจทย์รวมท้ายบท: ระดับ สอวน. ค่าย 1}

\subsection*{ระดับปรับพื้นฐาน}

\begin{enumerate}
    \item ภาชนะปริมาตร $V$ บรรจุแก๊สอุดมคติจำนวน $N$ โมเลกุล ที่อุณหภูมิ $T$ จงหาความดัน

    \item แก๊สปริมาณคงตัวอยู่ในถังแข็งที่อุณหภูมิ $T_0$ และความดัน $P_0$ ถ้าอุณหภูมิเพิ่มเป็น $3T_0/2$ จงหาความดันใหม่

    \item แก๊สอุดมคติอะตอมเดี่ยวจำนวน $n$ โมลมีอุณหภูมิเพิ่มขึ้น $\Delta T$ จงหาการเปลี่ยนแปลงพลังงานภายใน
\end{enumerate}

\subsection*{ระดับกลาง}

\begin{enumerate}
    \setcounter{enumi}{3}
    \item แก๊สในบอลลูนปริมาตรคงตัวมีความดันเกจ $P_g$ ในบริเวณที่ความดันบรรยากาศเป็น $P_a$ และอุณหภูมิ $T_1$ เมื่อบอลลูนลอยไปยังบริเวณที่อุณหภูมิ $T_2$ และความดันบรรยากาศ $P_b$ จงหาความดันเกจใหม่

    \item แก๊สอุดมคติอยู่ที่อุณหภูมิคงตัว $T$ ถ้าปริมาตรเปลี่ยนจาก $V$ เป็น $V+\Delta V$ โดย $|\Delta V|\ll V$ จงหา $\Delta P/P$ โดยประมาณ

    \item แก๊สอุดมคติอะตอมเดี่ยวจำนวน $n$ โมลอยู่ในกระบอกสูบที่มีลูกสูบลื่นและรับความร้อน $Q$ อย่างช้า ๆ ภายใต้ความดันภายนอกคงตัว จงหาอุณหภูมิที่เพิ่มขึ้น

    \item แก๊สอุดมคติสองชนิดมีอุณหภูมิเท่ากัน แต่โมเลกุลชนิดแรกมีมวล $m$ และชนิดที่สองมีมวล $4m$ จงหาอัตราส่วนพลังงานจลน์เฉลี่ย และอัตราส่วน $v_{\text{rms}}$
\end{enumerate}

\subsection*{ระดับยาก}

\begin{enumerate}
    \setcounter{enumi}{7}
    \item บอลลูนอากาศร้อนมีปริมาตร $V$ ภายในมีอากาศที่อุณหภูมิ $T_1$ ส่วนอากาศภายนอกมีอุณหภูมิ $T_0$ โดยความดันทั้งภายในและภายนอกเท่ากับ $P$ ให้มวลโมลาร์ของอากาศเป็น $M_{\text{mol}}$ จงหาแรงยกลัพธ์โดยไม่คิดมวลโครงสร้างบอลลูน

    \item แก๊สอุดมคติอะตอมเดี่ยวจำนวน $n$ โมลเปลี่ยนสถานะตามเส้นตรงบนกราฟ $P$ กับ $V$ จาก $(V_0,P_0)$ ไปยัง $(3V_0,2P_0)$ จงหางานที่แก๊สทำ และหา $\Delta U$ ในรูปของ $P_0V_0$

    \item ที่ความดันและอุณหภูมิเดียวกัน แก๊ส $A$ มีมวลโมลาร์ $M_A$ และแก๊ส $B$ มีมวลโมลาร์ $M_B$ จงหาอัตราส่วนความหนาแน่น และอัตราส่วนอัตราเร็วรากกำลังสองเฉลี่ย

    \item ให้ฟังก์ชันการกระจายอัตราเร็วเป็น
    \[
    N(v)=Av^2e^{-Bv^2/T}
    \]
    จงหา $v_{\text{peak}}$ ซึ่งทำให้ $N(v)$ มีค่าสูงสุด โดยไม่ต้องหาค่า $A$
\end{enumerate}

% =========================================================
\section{เฉลยย่อท้ายบท}

\begin{enumerate}
    \item
    \[
    \boxed{P=\frac{Nk_BT}{V}}
    \]

    \item
    \[
    \frac{P_0}{T_0}
    =
    \frac{P'}{3T_0/2}
    \]
    ดังนั้น
    \[
    \boxed{P'=\frac{3}{2}P_0}
    \]

    \item
    \[
    \boxed{\Delta U=\frac{3}{2}nR\Delta T}
    \]

    \item ความดันสัมบูรณ์เริ่มต้นคือ
    \[
    P_1=P_a+P_g
    \]
    เนื่องจากปริมาตรคงตัว
    \[
    \frac{P_1}{T_1}
    =
    \frac{P_2}{T_2}
    \]
    จึงได้
    \[
    P_2=(P_a+P_g)\frac{T_2}{T_1}
    \]
    ดังนั้น
    \[
    \boxed{
    P_{g,2}
    =
    (P_a+P_g)\frac{T_2}{T_1}-P_b
    }
    \]

    \item เพราะอุณหภูมิคงตัว
    \[
    PV=\text{ค่าคงตัว}
    \]
    ดังนั้น
    \[
    P\Delta V+V\Delta P\approx0
    \]
    \[
    \boxed{
    \frac{\Delta P}{P}
    \approx
    -\frac{\Delta V}{V}
    }
    \]

    \item
    \[
    Q
    =
    \Delta U+W
    \]
    \[
    Q
    =
    \frac{3}{2}nR\Delta T+nR\Delta T
    \]
    ดังนั้น
    \[
    \boxed{
    \Delta T
    =
    \frac{2Q}{5nR}
    }
    \]

    \item ที่อุณหภูมิเดียวกัน
    \[
    \left\langle K\right\rangle
    =
    \frac{3}{2}k_BT
    \]
    จึงเท่ากัน
    \[
    \boxed{
    \frac{\left\langle K_1\right\rangle}
    {\left\langle K_2\right\rangle}=1
    }
    \]
    ส่วน
    \[
    v_{\text{rms}}
    \propto
    \frac{1}{\sqrt{m_0}}
    \]
    จึงได้
    \[
    \boxed{
    \frac{v_{\text{rms},1}}
    {v_{\text{rms},2}}=2
    }
    \]

    \item
    \[
    \rho_{\text{out}}
    =
    \frac{PM_{\text{mol}}}{RT_0}
    \]
    \[
    \rho_{\text{in}}
    =
    \frac{PM_{\text{mol}}}{RT_1}
    \]
    ดังนั้น
    \[
    F_{\text{net}}
    =
    (\rho_{\text{out}}-\rho_{\text{in}})Vg
    \]
    \[
    \boxed{
    F_{\text{net}}
    =
    \frac{PM_{\text{mol}}Vg}{R}
    \left(
    \frac{1}{T_0}-\frac{1}{T_1}
    \right)
    }
    \]

    \item งานเท่ากับพื้นที่ใต้เส้นตรง
    \[
    W
    =
    \frac{P_0+2P_0}{2}(3V_0-V_0)
    \]
    \[
    \boxed{W=3P_0V_0}
    \]
    จากกฎแก๊ส
    \[
    nRT_i=P_0V_0
    \]
    \[
    nRT_f=(2P_0)(3V_0)=6P_0V_0
    \]
    ดังนั้นสำหรับแก๊สอะตอมเดี่ยว
    \[
    \Delta U
    =
    \frac{3}{2}nR(T_f-T_i)
    \]
    \[
    \boxed{
    \Delta U
    =
    \frac{15}{2}P_0V_0
    }
    \]

    \item
    \[
    \rho
    =
    \frac{PM_{\text{mol}}}{RT}
    \]
    ดังนั้น
    \[
    \boxed{
    \frac{\rho_A}{\rho_B}
    =
    \frac{M_A}{M_B}
    }
    \]
    และ
    \[
    v_{\text{rms}}
    =
    \sqrt{\frac{3RT}{M_{\text{mol}}}}
    \]
    จึงได้
    \[
    \boxed{
    \frac{v_{\text{rms},A}}
    {v_{\text{rms},B}}
    =
    \sqrt{\frac{M_B}{M_A}}
    }
    \]

    \item
    \[
    \frac{dN}{dv}
    =
    2Ave^{-Bv^2/T}
    \left(1-\frac{Bv^2}{T}\right)
    \]
    สำหรับยอดกราฟที่ $v>0$
    \[
    1-\frac{Bv^2}{T}=0
    \]
    ดังนั้น
    \[
    \boxed{
    v_{\text{peak}}
    =
    \sqrt{\frac{T}{B}}
    }
    \]
\end{enumerate}

% =========================================================
\section{สรุปสูตรสำคัญ}

\begin{tcolorbox}[title=สรุปบท: แก๊สอุดมคติ]
อุณหภูมิสัมบูรณ์

\[
T_{\si{K}}=T_{\si{\degreeCelsius}}+273.15
\]

กฎของแก๊สอุดมคติ

\[
PV=nRT
\]

\[
PV=Nk_BT
\]

\[
R=N_Ak_B
\]

แก๊สปริมาณคงตัว

\[
\frac{P_1V_1}{T_1}
=
\frac{P_2V_2}{T_2}
\]

ความดันสัมบูรณ์

\[
P_{\text{abs}}
=
P_{\text{atm}}+P_{\text{gauge}}
\]

ความหนาแน่นของแก๊ส

\[
\rho
=
\frac{PM_{\text{mol}}}{RT}
\]
\end{tcolorbox}

\begin{tcolorbox}[title=สรุปบท: ทฤษฎีจลน์และกฎข้อที่หนึ่ง]
ความดันจากการเคลื่อนที่ของโมเลกุล

\[
P
=
\frac{1}{3}
\left(\frac{N}{V}\right)
 m_0\overline{v^2}
\]

พลังงานจลน์เฉลี่ยของการเลื่อนที่

\[
\frac{1}{2}m_0\overline{v^2}
=
\frac{3}{2}k_BT
\]

อัตราเร็วรากกำลังสองเฉลี่ย

\[
 v_{\text{rms}}
=
\sqrt{\frac{3k_BT}{m_0}}
=
\sqrt{\frac{3RT}{M_{\text{mol}}}}
\]

หลักการแบ่งเท่าเทียมของพลังงาน

\[
\frac{1}{2}k_BT
\quad
\text{ต่อหนึ่ง degree of freedom}
\]

แก๊สอะตอมเดี่ยว

\[
U
=
\frac{3}{2}nRT
\]

\[
\Delta U
=
\frac{3}{2}nR\Delta T
\]

งานที่แก๊สทำ

\[
W
=
\int P\,dV
\]

กฎข้อที่หนึ่ง

\[
\Delta U
=
Q-W
\]
\end{tcolorbox}

% =========================================================
\section{ข้อสอบจริง สอวน. ที่ใช้ความรู้บทแก๊สอุดมคติและอุณหพลศาสตร์}


% =========================================================
\subsection{ปี 2560 ข้อ 19: จำนวนโมเลกุลของแก๊สไนโตรเจนในภาชนะ}

\vspace{1em}

\IfFileExists{img/posn60q19.png}{
\begin{center}
    \includegraphics[width=1\linewidth]{img/posn60q19.png}
\end{center}
}{
\begin{tcolorbox}[title=ตำแหน่งภาพที่แนะนำ]
ให้ตัดภาพข้อสอบปี 2560 ข้อ 19 แล้วบันทึกเป็น

\[
\texttt{img/posn60q19.png}
\]
\end{tcolorbox}
}

\begin{examplebox}{เฉลยละเอียด}
ภาชนะมีปริมาตร

\[
V=(0.25)(0.60)(0.50)
\]

\[
V=7.5\times10^{-2}\,\si{m^3}
\]

แก๊สมีอุณหภูมิ

\[
T=27+273=300\,\si{K}
\]

ความดันคือ

\[
P=0.05(1.013\times10^5)
\]

\[
P\approx5.1\times10^3\,\si{Pa}
\]

ใช้

\[
PV=nRT
\]

\[
n=\frac{PV}{RT}
\]

\[
n
=
\frac{(5.1\times10^3)(7.5\times10^{-2})}
{(8.31)(300)}
\]

\[
n\approx1.5\times10^{-1}\,\si{mol}
\]

จำนวนโมเลกุลคือ

\[
N=nN_A
\]

\[
N
\approx
(1.5\times10^{-1})(6.02\times10^{23})
\]

\[
N\approx9.2\times10^{22}
\]

ดังนั้นค่าประมาณตามลำดับขนาดคือ

\[
\boxed{N\approx10^{23}\text{ โมเลกุล}}
\]

\[
\boxed{\text{ตอบข้อ B}}
\]
\end{examplebox}

% =========================================================
\subsection{ปี 2561 ข้อ 17: ความดันเกจของแก๊สฮีเลียมในบอลลูน}

\vspace{1em}

\IfFileExists{img/posn61q17.png}{
\begin{center}
    \includegraphics[width=1\linewidth]{img/posn61q17.png}
\end{center}
}{
\begin{tcolorbox}[title=ตำแหน่งภาพที่แนะนำ]
ให้ตัดภาพข้อสอบปี 2561 ข้อ 17 แล้วบันทึกเป็น

\[
\texttt{img/posn61q17.png}
\]
\end{tcolorbox}
}

\begin{examplebox}{เฉลยละเอียด}
บอลลูนมีปริมาตรคงตัว ดังนั้นสำหรับแก๊สปริมาณเดิม

\[
\frac{P_1}{T_1}
=
\frac{P_2}{T_2}
\]

ต้องใช้ความดันสัมบูรณ์

ที่พื้นดิน

\[
P_1
=
P_{\text{atm},1}+P_{\text{gauge},1}
\]

\[
P_1=100+260
\]

\[
P_1=360\,\si{kPa}
\]

อุณหภูมิเริ่มต้นคือ

\[
T_1=27+273=300\,\si{K}
\]

ที่ระดับความสูง $5.0\,\si{km}$

\[
T_2=-23+273=250\,\si{K}
\]

ดังนั้น

\[
P_2
=
P_1\frac{T_2}{T_1}
\]

\[
P_2
=
360\left(\frac{250}{300}\right)
\]

\[
P_2=300\,\si{kPa}
\]

นี่คือความดันสัมบูรณ์ ความดันเกจใหม่จึงเป็น

\[
P_{\text{gauge},2}
=
P_2-P_{\text{atm},2}
\]

\[
P_{\text{gauge},2}=300-60
\]

\[
\boxed{P_{\text{gauge},2}=240\,\si{kPa}}
\]

\[
\boxed{\text{ตอบข้อ B}}
\]
\end{examplebox}

% =========================================================
\subsection{ปี 2561 ข้อ 19: แก๊สในกระป๋องและแรงยกฝา}

\vspace{1em}

\IfFileExists{img/posn61q19.png}{
\begin{center}
    \includegraphics[width=1\linewidth]{img/posn61q19.png}
\end{center}
}{
\begin{tcolorbox}[title=ตำแหน่งภาพที่แนะนำ]
ให้ตัดภาพข้อสอบปี 2561 ข้อ 19 แล้วบันทึกเป็น

\[
\texttt{img/posn61q19.png}
\]
\end{tcolorbox}
}

\begin{examplebox}{เฉลยละเอียด}
กระป๋องทรงกระบอกมีความสูง

\[
h=20\,\si{cm}=0.20\,\si{m}
\]

ฝากระป๋องมีพื้นที่

\[
A=40\,\si{cm^2}=4.0\times10^{-3}\,\si{m^2}
\]

ภายในบรรจุแก๊ส

\[
n=\frac{1}{83}\,\si{mol}
\]

ให้ $P_g$ เป็นความดันสัมบูรณ์ของแก๊สภายใน

เมื่อเริ่มยกฝาได้พอดี แรงขึ้นเท่ากับแรงลง

\[
F+P_gA=P_{\text{atm}}A
\]

ดังนั้น

\[
P_g
=
P_{\text{atm}}-\frac{F}{A}
\]

แทนค่า

\[
P_g
=
100\times10^3
-
\frac{240}{4.0\times10^{-3}}
\]

\[
P_g=40\times10^3\,\si{Pa}
\]

ปริมาตรแก๊สคือ

\[
V=Ah
\]

\[
V
=
(4.0\times10^{-3})(0.20)
\]

\[
V=8.0\times10^{-4}\,\si{m^3}
\]

ใช้

\[
PV=nRT
\]

\[
T=\frac{PV}{nR}
\]

แทนค่า

\[
T
=
\frac{(40\times10^3)(8.0\times10^{-4})}
{(1/83)(8.3)}
\]

\[
\boxed{T=320\,\si{K}}
\]

\[
\boxed{\text{ตอบข้อ A}}
\]
\end{examplebox}

% =========================================================
\subsection{ปี 2562 ข้อ 16: ให้ความร้อนแก่แก๊สอะตอมเดี่ยวในกระบอกสูบ}

\vspace{1em}

\IfFileExists{img/posn62q16.png}{
\begin{center}
    \includegraphics[width=1\linewidth]{img/posn62q16.png}
\end{center}
}{
\begin{tcolorbox}[title=ตำแหน่งภาพที่แนะนำ]
ให้ตัดภาพข้อสอบปี 2562 ข้อ 16 แล้วบันทึกเป็น

\[
\texttt{img/posn62q16.png}
\]
\end{tcolorbox}
}

\begin{examplebox}{เฉลยละเอียด}
แก๊สอุดมคติอะตอมเดี่ยวจำนวน $n$ โมล อยู่ในกระบอกสูบที่มีลูกสูบลื่น และให้ความร้อนอย่างช้า ๆ ปริมาณ $Q$

เพราะลูกสูบลื่นและภายนอกมีความดันบรรยากาศคงตัว กระบวนการจึงเป็นกระบวนการความดันคงตัว

สำหรับแก๊สอะตอมเดี่ยว

\[
\Delta U
=
\frac{3}{2}nR\Delta T
\]

งานที่แก๊สทำคือ

\[
W=P\Delta V
\]

จากกฎแก๊สอุดมคติที่ความดันคงตัว

\[
P\Delta V=nR\Delta T
\]

ดังนั้น

\[
W=nR\Delta T
\]

ใช้กฎข้อที่หนึ่ง

\[
Q=\Delta U+W
\]

\[
Q
=
\frac{3}{2}nR\Delta T+nR\Delta T
\]

\[
Q
=
\frac{5}{2}nR\Delta T
\]

ดังนั้น

\[
\boxed{
T_1-T_0
=
\Delta T
=
\frac{2Q}{5nR}
}
\]
\end{examplebox}

% =========================================================
\subsection{อ่านเสริม: ข้อสอบศูนย์ สอวน. ปี 2563 ข้อ 10 ความหนาแน่นของอากาศและไฮโดรเจน}

\vspace{1em}

\IfFileExists{img/posn63centerq10.png}{
\begin{center}
    \includegraphics[width=1\linewidth]{img/posn63centerq10.png}
\end{center}
}{
\begin{tcolorbox}[title=ตำแหน่งภาพที่แนะนำ]
ให้ตัดภาพข้อสอบศูนย์ สอวน. ปี 2563 ข้อ 10 แล้วบันทึกเป็น

\[
\texttt{img/posn63centerq10.png}
\]
\end{tcolorbox}
}

\begin{examplebox}{เฉลยละเอียด}
สำหรับแก๊สอุดมคติ

\[
\rho
=
\frac{PM_{\text{mol}}}{RT}
\]

อากาศและแก๊สไฮโดรเจนอยู่ที่ความดันและอุณหภูมิเดียวกัน จึงตัด $P$, $R$, และ $T$ ออกได้เมื่อหาอัตราส่วน

\[
\frac{\rho_{\text{air}}}{\rho_{H_2}}
=
\frac{PM_{\text{air}}/(RT)}{PM_H/(RT)}
\]

ดังนั้น

\[
\boxed{
\frac{\rho_{\text{air}}}{\rho_{H_2}}
=
\frac{M_{\text{air}}}{M_H}
}
\]
\end{examplebox}

% =========================================================
\subsection{อ่านเสริม: ข้อสอบศูนย์ สอวน. ปี 2563 ข้อ 11 น้ำแข็งแห้งกลายเป็นแก๊ส}

\vspace{1em}

\IfFileExists{img/posn63centerq11.png}{
\begin{center}
    \includegraphics[width=1\linewidth]{img/posn63centerq11.png}
\end{center}
}{
\begin{tcolorbox}[title=ตำแหน่งภาพที่แนะนำ]
ให้ตัดภาพข้อสอบศูนย์ สอวน. ปี 2563 ข้อ 11 แล้วบันทึกเป็น

\[
\texttt{img/posn63centerq11.png}
\]
\end{tcolorbox}
}

\begin{examplebox}{เฉลยละเอียด}
น้ำแข็งแห้งมีปริมาตรเริ่มต้น

\[
V_s=1.0\,\si{cm^3}
\]

และความหนาแน่น

\[
\rho_s=1.56\,\si{g/cm^3}
\]

ดังนั้นมวลคือ

\[
m=\rho_sV_s
\]

\[
m=1.56\,\si{g}
\]

มวลโมลาร์ของแก๊สคาร์บอนไดออกไซด์คือ

\[
M_{\text{CO}_2}=12+2(16)=44\,\si{g/mol}
\]

จำนวนโมลคือ

\[
n=\frac{m}{M_{\text{CO}_2}}
\]

\[
n=\frac{1.56}{44}\,\si{mol}
\]

ที่

\[
T=0+273=273\,\si{K}
\]

และ

\[
P=1.01325\times10^5\,\si{Pa}
\]

ใช้

\[
PV=nRT
\]

\[
V=\frac{nRT}{P}
\]

\[
V
=
\frac{(1.56/44)(8.314)(273)}{1.01325\times10^5}
\]

\[
V\approx7.94\times10^{-4}\,\si{m^3}
\]

เปลี่ยนหน่วย

\[
1\,\si{m^3}=10^6\,\si{cm^3}
\]

ดังนั้น

\[
\boxed{
V\approx7.94\times10^2\,\si{cm^3}
}
\]
\end{examplebox}

% =========================================================
\subsection{อ่านเสริม: ข้อสอบศูนย์ สอวน. ปี 2563 ข้อ 29 แรงยกของบอลลูนอากาศร้อน}

\vspace{1em}

\IfFileExists{img/posn63centerq29.png}{
\begin{center}
    \includegraphics[width=1\linewidth]{img/posn63centerq29.png}
\end{center}
}{
\begin{tcolorbox}[title=ตำแหน่งภาพที่แนะนำ]
ให้ตัดภาพข้อสอบศูนย์ สอวน. ปี 2563 ข้อ 29 แล้วบันทึกเป็น

\[
\texttt{img/posn63centerq29.png}
\]
\end{tcolorbox}
}

\begin{examplebox}{เฉลยละเอียด}
ให้บอลลูนมีปริมาตร

\[
V
\]

ความดันอากาศภายในและภายนอกเท่ากันคือ

\[
P
\]

อากาศภายนอกมีอุณหภูมิ

\[
T_0
\]

และอากาศภายในมีอุณหภูมิ

\[
T_1
\]

โดยทั่วไป

\[
T_1>T_0
\]

ให้มวลโมลาร์ของอากาศเป็น

\[
M
\]

ความหนาแน่นของอากาศภายนอกคือ

\[
\rho_{\text{out}}
=
\frac{PM}{RT_0}
\]

ความหนาแน่นของอากาศภายในคือ

\[
\rho_{\text{in}}
=
\frac{PM}{RT_1}
\]

แรงพยุงที่กระทำต่อบอลลูนเท่ากับน้ำหนักของอากาศภายนอกที่ถูกแทนที่

\[
B=\rho_{\text{out}}Vg
\]

น้ำหนักของอากาศร้อนภายในคือ

\[
W_{\text{in}}=\rho_{\text{in}}Vg
\]

ดังนั้นแรงยกลัพธ์ เมื่อไม่คิดมวลโครงสร้างบอลลูน คือ

\[
F_{\text{lift}}
=
B-W_{\text{in}}
\]

\[
F_{\text{lift}}
=
(\rho_{\text{out}}-\rho_{\text{in}})Vg
\]

จึงได้

\[
\boxed{
F_{\text{lift}}
=
\frac{PMVg}{R}
\left(
\frac{1}{T_0}
-
\frac{1}{T_1}
\right)
}
\]
\end{examplebox}

% =========================================================
\subsection{ปี 2564 ข้อ 14: การเปลี่ยนแปลงเล็กน้อยของความดันและปริมาตรที่อุณหภูมิคงตัว}

\vspace{1em}

\IfFileExists{img/posn64q14.png}{
\begin{center}
    \includegraphics[width=1\linewidth]{img/posn64q14.png}
\end{center}
}{
\begin{tcolorbox}[title=ตำแหน่งภาพที่แนะนำ]
ให้ตัดภาพข้อสอบปี 2564 ข้อ 14 แล้วบันทึกเป็น

\[
\texttt{img/posn64q14.png}
\]
\end{tcolorbox}
}

\begin{examplebox}{เฉลยละเอียด}
แก๊สอุดมคติมีอุณหภูมิคงตัว

\[
PV=nRT=\text{ค่าคงตัว}
\]

เมื่อความดันและปริมาตรเปลี่ยนไปเล็กน้อย

\[
(P+\Delta P)(V+\Delta V)=PV
\]

กระจายพจน์

\[
PV+P\Delta V+V\Delta P+\Delta P\Delta V=PV
\]

เพราะการเปลี่ยนแปลงมีขนาดเล็ก ละพจน์อันดับสอง

\[
\Delta P\Delta V
\]

จะได้

\[
P\Delta V+V\Delta P\approx0
\]

ดังนั้น

\[
\boxed{
\frac{\Delta P}{\Delta V}
=
-\frac{P}{V}
}
\]

\[
\boxed{\text{ตอบข้อ ข.}}
\]
\end{examplebox}

% =========================================================
\subsection{ปี 2565 ข้อ 14: น้ำกลายเป็นไอน้ำที่ความดันคงตัว}

\vspace{1em}

\IfFileExists{img/posn65q14.png}{
\begin{center}
    \includegraphics[width=1\linewidth]{img/posn65q14.png}
\end{center}
}{
\begin{tcolorbox}[title=ตำแหน่งภาพที่แนะนำ]
ให้ตัดภาพข้อสอบปี 2565 ข้อ 14 แล้วบันทึกเป็น

\[
\texttt{img/posn65q14.png}
\]
\end{tcolorbox}
}

\begin{examplebox}{เฉลยละเอียด}
น้ำมวล

\[
m=1.0\,\si{g}
\]

เมื่อกลายเป็นไอน้ำหมดที่

\[
T=100+273=373\,\si{K}
\]

จำนวนโมลของไอน้ำคือ

\[
n=\frac{1}{18}\,\si{mol}
\]

ใช้กฎของแก๊สอุดมคติ

\[
PV=nRT
\]

\[
V_{\text{steam}}
=
\frac{nRT}{P}
\]

\[
V_{\text{steam}}
=
\left(\frac{373}{18}\right)
\left(\frac{R}{P}\right)
\si{m^3}
\]

น้ำเหลวมวล $1\,\si{g}$ มีปริมาตรประมาณ

\[
V_{\text{water}}=1\,\si{cm^3}=10^{-6}\,\si{m^3}
\]

ดังนั้นอัตราส่วนปริมาตรคือ

\[
\frac{V_{\text{steam}}}{V_{\text{water}}}
=
\left(\frac{373}{18}\right)
\left(\frac{R}{P}\right)
\times10^6
\]

\[
\boxed{
\frac{V_{\text{steam}}}{V_{\text{water}}}
=
\left(\frac{373}{18}\right)
\left(\frac{R}{P}\right)
\times10^6
}
\]

\[
\boxed{\text{ตอบข้อ ข.}}
\]
\end{examplebox}

% =========================================================
\subsection{ปี 2566 ข้อ 5: กราฟการกระจายอัตราเร็วของโมเลกุล}

\vspace{1em}

\IfFileExists{img/posn66q5.png}{
\begin{center}
    \includegraphics[width=1\linewidth]{img/posn66q5.png}
\end{center}
}{
\begin{tcolorbox}[title=ตำแหน่งภาพที่แนะนำ]
ให้ตัดภาพข้อสอบปี 2566 ข้อ 5 แล้วบันทึกเป็น

\[
\texttt{img/posn66q5.png}
\]
\end{tcolorbox}
}

\begin{examplebox}{เฉลยละเอียด}
โจทย์ให้

\[
N(v)
=
Av^2e^{-Bv^2/T}
\]

เมื่อ

\[
v=0
\]

จะได้

\[
N(0)=0
\]

ดังนั้นกราฟต้องเริ่มจากจุดกำเนิด

สำหรับ $v$ ขนาดเล็ก

\[
e^{-Bv^2/T}\approx1
\]

จึงมี

\[
N(v)\approx Av^2
\]

กราฟจึงเริ่มเพิ่มขึ้นอย่างราบเรียบ และมีความชันเป็นศูนย์ที่จุดกำเนิด

เมื่อ $v$ มีค่ามาก

\[
e^{-Bv^2/T}\longrightarrow0
\]

ดังนั้น

\[
N(v)\longrightarrow0
\]

กราฟที่ถูกต้องจึงต้องเริ่มจากศูนย์ โค้งขึ้นอย่างราบเรียบ มีค่าสูงสุดหนึ่งครั้ง แล้วลดลงเข้าใกล้ศูนย์

\[
\boxed{\text{ตอบข้อ 4}}
\]
\end{examplebox}

% =========================================================
\subsection{โจทย์เชื่อมโยงกับของไหล: ปี 2566 ข้อ 7 จำนวนโมเลกุลทั้งหมดในบรรยากาศโลก}

\vspace{1em}

\IfFileExists{img/posn66q7.png}{
\begin{center}
    \includegraphics[width=1\linewidth]{img/posn66q7.png}
\end{center}
}{
\begin{tcolorbox}[title=ตำแหน่งภาพที่แนะนำ]
ให้ตัดภาพข้อสอบปี 2566 ข้อ 7 แล้วบันทึกเป็น

\[
\texttt{img/posn66q7.png}
\]
\end{tcolorbox}
}

\begin{examplebox}{เฉลยละเอียด}
ความดันบรรยากาศที่ผิวโลกเกิดจากน้ำหนักรวมของบรรยากาศต่อพื้นที่ผิวโลก

พื้นที่ผิวโลกคือ

\[
A=4\pi R^2
\]

ถ้ามวลรวมของบรรยากาศเป็น $M_{\text{atm}}$

\[
P
=
\frac{M_{\text{atm}}g}{4\pi R^2}
\]

ดังนั้น

\[
M_{\text{atm}}
=
\frac{4\pi R^2P}{g}
\]

ให้มวลของโมเลกุลอากาศหนึ่งโมเลกุลเป็น

\[
m_0\approx5\times10^{-26}\,\si{kg}
\]

จำนวนโมเลกุลทั้งหมดคือ

\[
N
=
\frac{M_{\text{atm}}}{m_0}
\]

\[
N
=
\frac{4\pi R^2P}{m_0g}
\]

แทนค่าโดยประมาณ

\[
N
\approx
\frac{4\pi(6\times10^6)^2(1\times10^5)}
{(5\times10^{-26})(10)}
\]

\[
N\approx9\times10^{43}
\]

ดังนั้นลำดับขนาดคือ

\[
\boxed{N\approx10^{44}\text{ โมเลกุล}}
\]

\[
\boxed{\text{ตอบข้อ 2}}
\]
\end{examplebox}

% =========================================================
\subsection{ปี 2566 ข้อ 24: น้ำกลายเป็นไอน้ำที่อุณหภูมิ $267\,\si{\degreeCelsius}$}

\vspace{1em}

\IfFileExists{img/posn66q24.png}{
\begin{center}
    \includegraphics[width=1\linewidth]{img/posn66q24.png}
\end{center}
}{
\begin{tcolorbox}[title=ตำแหน่งภาพที่แนะนำ]
ให้ตัดภาพข้อสอบปี 2566 ข้อ 24 แล้วบันทึกเป็น

\[
\texttt{img/posn66q24.png}
\]
\end{tcolorbox}
}

\begin{examplebox}{เฉลยละเอียด}
น้ำเหลวเริ่มต้นมีปริมาตร

\[
1\,\si{cm^3}
\]

จึงมีมวลประมาณ

\[
1\,\si{g}
\]

จำนวนโมลของน้ำคือ

\[
n=\frac{1}{18}\,\si{mol}
\]

อุณหภูมิของไอน้ำคือ

\[
T=267+273
\]

\[
T=540\,\si{K}
\]

ความดันคือ

\[
P=1.0\times10^5\,\si{Pa}
\]

ใช้กฎของแก๊สอุดมคติ

\[
PV=nRT
\]

\[
V=\frac{nRT}{P}
\]

\[
V
=
\frac{(1/18)(8.3)(540)}{1.0\times10^5}
\]

\[
V=2.49\times10^{-3}\,\si{m^3}
\]

เปลี่ยนหน่วย

\[
1\,\si{m^3}=10^6\,\si{cm^3}
\]

จึงได้

\[
\boxed{
V\approx2.49\times10^3\,\si{cm^3}
}
\]
\end{examplebox}